\documentclass{article}
\usepackage[en]{homework}

\config{Algebra\ -\ 0}{2026 Spring}{1}

\begin{document}

\maketitle

\begin{problem}[1B-2]
Suppose $ a\in\mathbb{F} $, $ v\in V $, and $ av=0 $. Prove that: $ a=0 $ or $ v=0 $.
\end{problem}

\begin{proof}
Assume that $ a\neq 0 $, then it's sufficient to prove $ v=0 $.\par 
By the associativity of the scalar multiplication, $ v=(a^{-1}a)v=a^{-1}(av)=a^{-1}0 $. By the property of the zero vector, $ a^{-1}0=0 $. Thus $ v=0 $.
\end{proof}

\begin{problem}[1B-3]
Suppose $v, w \in V$. Explain why there exists a unique $x \in V$ such that $v + 3x = w$.
\end{problem}

\begin{solution}
The addictive inverse of $ v $ is unique, denoted as $ -v $. Then $ x=\frac{1}{3}(w-v) $ is the unique solution to the equation $ v+3x=w $.
\end{solution}

\begin{problem}[1B-5]
Show that in the definition of a vector space (1.20), the additive inverse condition can be replaced with the condition that
\[ 0v=0 \text{ for all }v\in V. \]
Here the 0 on the left side is the number 0, and the 0 on the right side is the additive identity of $ V $.
\end{problem}

\begin{proof}
For every $ v\in V $, define the addictive inverse of $ v $ is $ (-1)v $. Then $ v+(-1)v=v+(-1)1v=(1-1)v=0v=0 $. Thus the addictive inverse condition is satisfied.
\end{proof}

\begin{problem}[1B-6]
Let $\infty$ and $-\infty$ denote two distinct objects, neither of which is in $\R$. Define an addition and scalar multiplication on $\R \cup \{ \infty, -\infty \}$ as you could guess from the notation. Specifically, the sum and product of two real numbers is as usual, and for $t \in \R$ define
\[
t\infty = 
\begin{cases} 
-\infty & \text{if } t < 0, \\
0 & \text{if } t = 0, \\
\infty & \text{if } t > 0.
\end{cases}
\]
\[
t(-\infty) = 
\begin{cases} 
\infty & \text{if } t < 0, \\
0 & \text{if } t = 0, \\
-\infty & \text{if } t > 0.
\end{cases}
\]
and
\[
t + \infty = \infty + t = \infty + \infty = \infty,
\]
\[
t + (-\infty) = (-\infty) + t = (-\infty) + (-\infty) = -\infty,
\]
\[
\infty + (-\infty) = (-\infty) + \infty = 0.
\]
With these operations of addition and scalar multiplication, is $\R \cup \{ \infty, -\infty \}$ a vector space over $\R$? Explain.
\end{problem}

\begin{solution}
No. The distributive law fails, since
\[ (-\infty+\infty)+\infty= 0+\infty = \infty,\text{ but } -\infty+(\infty+\infty) = -\infty+\infty = 0. \]
\end{solution}

\begin{problem}[1C-1]
For each of the following subsets of $\mathbb{F}^3$, determine whether it is a subspace of $\mathbb{F}^3$.

\begin{enumerate}[label=(\alph*)]
    \item $\{ (x_1, x_2, x_3) \in \mathbb{F}^3 : x_1 + 2x_2 + 3x_3 = 0 \}$
    \item $\{ (x_1, x_2, x_3) \in \mathbb{F}^3 : x_1 + 2x_2 + 3x_3 = 4 \}$
    \item $\{ (x_1, x_2, x_3) \in \mathbb{F}^3 : x_1 x_2 x_3 = 0 \}$
    \item $\{ (x_1, x_2, x_3) \in \mathbb{F}^3 : x_1 = 5x_3 \}$
\end{enumerate}
\end{problem}

\begin{solution}
\begin{enumerate}[label=(\alph*)]
    \item Yes. The zero vector is in the set, and the set is closed under addition and scalar multiplication.
    \item No. The zero vector is not in the set.
    \item No. The set is not closed under addition, since $ (1, 0, 0) $ and $ (0, 1, 1) $ are in the set but their sum $ (1, 1, 1) $ is not in the set.
    \item Yes. The reason is similar to (a).
\end{enumerate}
\end{solution}

\begin{problem}[1C-5]
Is $ \R^2 $ is a subspace of the complex space $ \C^2 $?
\end{problem}

\begin{solution}
No, since \(\R^2\) is not closed under scalar multiplication with respect to an imaginary scalar in \(\C\).
\end{solution}

\begin{problem}[1C-6]
\begin{enumerate}[label=(\alph*)]
    \item Is $\{(a, b, c) \in \R^3 : a^3 = b^3\}$ a subspace of $\R^3$?
    \item Is $\{(a, b, c) \in \C^3 : a^3 = b^3\}$ a subspace of $\C^3$?
\end{enumerate}
\end{problem}

\begin{solution}
\begin{enumerate}[label=(\alph*)]
    \item Yes. $ a^3=b^3, a,b\in\R $ implies $ a=b $, then the set is equal to $\{(a, b, c) \in \R^3 : a = b\}$, which is a subspace of $\R^3$.
    \item No. The set is not closed under addition, since $ (1, 1, 0) $ and $ (1, \omega, 0) $ are in the set but their sum $ (2, 1+\omega, 0) $ is not in the set, where $ \omega=(-1+\sqrt{3}\i)/2 $.
\end{enumerate}
\end{solution}

\begin{problem}[1C-12]
Prove that the union of two subspaces of $ V $ is a subspace of $ V $ if and only if one of the subspaces is contained in the other.
\end{problem}

\begin{proof}
The sufficiency is trivial, so we will only prove the necessity.\par 
Conversly, assume that $ U_1, U_2 $ and $ U:=U_1\cup U_2 $ are subspaces of $ V $, but neither $ U_1\subset U_2 $ nor $ U_2\subset U_1 $. Then, there exist $ u_1\in U_1 - U_2 $ and $ u_2\in U_2 - U_1 $.\par
Since $ U $ is closed under addition, $ u_1+u_2\in U $, which implise ether $ u_1+u_2\in U_1 $ or $ u_1+u_2\in U_2 $. However, if $ u_1+u_2\in U_1 $, then $ u_2=(u_1+u_2)-u_1\in U_1 $, which contradicts the choice of $ u_2 $. Similarly, if $ u_1+u_2\in U_2 $, then $ u_1\in U_2 $, which contradicts the choice of $ u_1 $. Thus we get a contradiction, and the necessity is proved. 
\end{proof}

\begin{problem}[1C-13]
Prove that the union of three subspaces of $ V $ is a subspace of $ V $ if and only if one of the subspaces contains the other two.
\end{problem}

\end{document}